% Generated draft: requires mathematical and source review.
\documentclass[11pt]{article}
\usepackage[margin=25mm]{geometry}
\usepackage{fontspec}
\setmainfont{Latin Modern Roman}
\usepackage{amsmath,amssymb,mathtools,unicode-math}
\setmathfont{Latin Modern Math}
\usepackage{enumitem,needspace,xurl,hyperref}
\setlength{\parindent}{0pt}
\setlength{\parskip}{5pt}
\setlength{\emergencystretch}{3em}
\newcommand{\N}{\mathbb{N}}
\newcommand{\Z}{\mathbb{Z}}
\newcommand{\Q}{\mathbb{Q}}
\newcommand{\R}{\mathbb{R}}
\newcommand{\C}{\mathbb{C}}
\newcommand{\F}{\mathbb{F}}
\newcommand{\A}{\mathbb{A}}
\newcommand{\PP}{\mathbb{P}}
\newcommand{\E}{\mathbb{E}}
\newcommand{\eps}{\varepsilon}
\newcommand{\dd}{\,\mathrm{d}}
\begin{document}
\section*{UnsolvedMath ID: 800353}
\textbf{Source problem code:} catalogue entry
\subsection{Definitions and mathematical statement}
Let \(S_n\) denote the large Schr\"oder numbers with the conventional normalization
\[
S_0=1,\qquad S_n=S_{n-1}+\sum_{k=0}^{n-1}S_kS_{n-1-k}\quad(n\in\Z,\ n\geq1).
\]
The supplied record uses the notation \(T_{112}(a,b)\) from Propp's cited source but does not define it. In particular, the counted objects, the underlying region, the parameter conventions, and the meaning of the subscript \(112\) are not supplied. These definitions must be recovered from the cited source before the question becomes self-contained.

Subject to that unresolved definition and confirmation that the source's phrase ``\(n\)th large Schr\"oder number'' uses the normalization above, the proposed statement is
\[
\forall n\in\Z,\quad n\geq1\ \Longrightarrow\ T_{112}(n+2,2n+1)=S_n.
\]
The source record reports that equality has been verified for \(1\leq n\leq15\); this is a supplied verification claim, not an independently checked result. The registry labels the problem open.
\subsection{Short English statement}
Propp asks whether \(T_{112}(n+2,2n+1)\) equals the \(n\)th large Schr\"oder number for every integer \(n\geq1\). The supplied record reports agreement through \(n=15\). A complete formulation requires the missing definition of \(T_{112}\) and confirmation of the sequence indexing.
\subsection{Sources}
\begin{itemize}
\item Local UnsolvedMath registry, permanent ID \(800353\), source code \texttt{catalogue entry}, title: \emph{Trimer Covers in the Triangular Grid --- Problem 12}. The supplied record attributes the problem to Propp (2022) and labels it open; that status is supplied curation, not an independently verified literature assessment.
\item Propp, \emph{Trimer Covers in the Triangular Grid: Twenty Mostly Open Problems} (2022), source item \texttt{Problem 12 (\#5)}. Supplied source URL: \url{https://arxiv.org/abs/2206.06472}. The record reports access on 2026-10-16 and extraction by \texttt{arxiv\_latex}; this provenance has not been independently verified here.
\item OEIS entry A006318, referenced in the supplied statement for the large Schr\"oder numbers: \url{https://oeis.org/A006318}.
\end{itemize}
\end{document}
